% Section pour intégrer toutes les captures d'écran dans le rapport AirOps

\section{Captures d'Écran et Démonstrations}

Cette section présente les interfaces et fonctionnalités réalisées de la solution AirOps à travers des captures d'écran détaillées de l'application en fonctionnement.

\subsection{Interface Utilisateur Frontend}

\subsubsection{Authentification et Sécurité}

\begin{figure}[H]
    \centering
    \includegraphics[width=0.8\textwidth]{screenshots/01-frontend/login.png}
    \caption{Page de connexion AirOps avec authentification JWT}
    \label{fig:airops-login}
\end{figure}

La page de connexion présente une interface moderne et sécurisée utilisant l'authentification JWT. Les utilisateurs peuvent se connecter avec leurs identifiants pour accéder aux fonctionnalités selon leur rôle (ADMIN, SUPERVISEUR, AGENT).

\subsubsection{Tableau de Bord Principal}

\begin{figure}[H]
    \centering
    \includegraphics[width=1.0\textwidth]{screenshots/01-frontend/dashboard.png}
    \caption{Tableau de bord principal avec métriques temps réel}
    \label{fig:airops-dashboard}
\end{figure}

Le tableau de bord offre une vue d'ensemble complète des opérations avec :
\begin{itemize}
    \item Statistiques des vols en temps réel
    \item Indicateurs de performance du personnel
    \item Notifications et alertes système
    \item Accès rapide aux fonctionnalités principales
\end{itemize}

\subsubsection{Gestion des Vols}

\begin{figure}[H]
    \centering
    \includegraphics[width=1.0\textwidth]{screenshots/01-frontend/flights-list.png}
    \caption{Interface de gestion des vols avec filtres et recherche}
    \label{fig:airops-flights-list}
\end{figure}

L'interface de gestion des vols permet :
\begin{itemize}
    \item Visualisation complète des vols planifiés
    \item Filtrage par statut, date, origine et destination
    \item Mise à jour temps réel des statuts de vol
    \item Pagination pour la gestion de volumes importants
\end{itemize}

\begin{figure}[H]
    \centering
    \includegraphics[width=0.9\textwidth]{screenshots/01-frontend/flight-form.png}
    \caption{Formulaire de création/modification de vol}
    \label{fig:airops-flight-form}
\end{figure}

Le formulaire de gestion des vols intègre :
\begin{itemize}
    \item Validation en temps réel des données saisies
    \item Sélecteurs d'aéroports avec autocomplétion
    \item Gestion des horaires avec validation de cohérence
    \item Interface responsive adaptée aux différents appareils
\end{itemize}

\subsubsection{Gestion du Personnel}

\begin{figure}[H]
    \centering
    \includegraphics[width=1.0\textwidth]{screenshots/01-frontend/staff-management.png}
    \caption{Interface de gestion du personnel avec profils complets}
    \label{fig:airops-staff-management}
\end{figure}

Le module RH offre une gestion complète du personnel avec :
\begin{itemize}
    \item Profils détaillés avec photos et qualifications
    \item Filtrage par département, rôle et statut
    \item Gestion des plannings et horaires de travail
    \item Système d'évaluation de performance intégré
\end{itemize}

\begin{figure}[H]
    \centering
    \includegraphics[width=0.9\textwidth]{screenshots/01-frontend/staff-profile.png}
    \caption{Profil détaillé d'un membre du personnel}
    \label{fig:airops-staff-profile}
\end{figure}

\subsubsection{Communication Temps Réel}

\begin{figure}[H]
    \centering
    \includegraphics[width=1.0\textwidth]{screenshots/01-frontend/chat-interface.png}
    \caption{Interface de communication temps réel avec WebSocket}
    \label{fig:airops-chat-interface}
\end{figure}

La plateforme de communication intègre :
\begin{itemize}
    \item Chat public pour communications générales
    \item Messages privés entre utilisateurs
    \item Salles de discussion thématiques
    \item Notifications push en temps réel
    \item Historique des conversations avec recherche
\end{itemize}

\subsubsection{Interface Responsive}

\begin{figure}[H]
    \centering
    \includegraphics[width=0.6\textwidth]{screenshots/01-frontend/mobile-dashboard.png}
    \caption{Version mobile du tableau de bord}
    \label{fig:airops-mobile}
\end{figure}

L'interface s'adapte automatiquement aux différentes tailles d'écran, garantissant une expérience utilisateur optimale sur tous les appareils.

\subsection{Architecture Backend et Services}

\subsubsection{Service Discovery - Eureka}

\begin{figure}[H]
    \centering
    \includegraphics[width=1.0\textwidth]{screenshots/02-backend/eureka-dashboard.png}
    \caption{Dashboard Eureka avec services enregistrés}
    \label{fig:eureka-dashboard}
\end{figure}

Le registre de services Eureka affiche tous les microservices actifs avec leur état de santé et leurs instances disponibles.

\subsubsection{API Gateway}

\begin{figure}[H]
    \centering
    \includegraphics[width=1.0\textwidth]{screenshots/02-backend/api-gateway-routes.png}
    \caption{Configuration des routes API Gateway}
    \label{fig:api-gateway-routes}
\end{figure}

\subsubsection{Health Checks et Monitoring}

\begin{figure}[H]
    \centering
    \includegraphics[width=0.8\textwidth]{screenshots/02-backend/actuator-health.png}
    \caption{Endpoints Spring Boot Actuator pour monitoring}
    \label{fig:actuator-health}
\end{figure}

Les endpoints Actuator fournissent des informations détaillées sur l'état de santé et les métriques de chaque service.

\subsection{Containerisation Docker}

\subsubsection{Containers en Fonctionnement}

\begin{figure}[H]
    \centering
    \includegraphics[width=1.0\textwidth]{screenshots/03-docker/containers-running.png}
    \caption{Containers Docker actifs avec docker-compose}
    \label{fig:docker-containers}
\end{figure}

L'architecture containerisée montre tous les services en fonctionnement :
\begin{itemize}
    \item 8 microservices Spring Boot
    \item 3 instances PostgreSQL (Auth, Flight, Staff)
    \item 1 instance MongoDB (Chat)
    \item Frontend Next.js
    \item Services d'infrastructure (Eureka, Gateway, Config)
\end{itemize}

\subsubsection{Images Docker}

\begin{figure}[H]
    \centering
    \includegraphics[width=1.0\textwidth]{screenshots/03-docker/docker-images.png}
    \caption{Images Docker construites pour les services AirOps}
    \label{fig:docker-images}
\end{figure}

\subsubsection{Monitoring des Ressources}

\begin{figure}[H]
    \centering
    \includegraphics[width=1.0\textwidth]{screenshots/03-docker/docker-stats.png}
    \caption{Utilisation des ressources par les containers}
    \label{fig:docker-stats}
\end{figure}

\subsection{Bases de Données}

\subsubsection{PostgreSQL - Données Relationnelles}

\begin{figure}[H]
    \centering
    \includegraphics[width=1.0\textwidth]{screenshots/04-databases/postgresql-auth-tables.png}
    \caption{Structure de la base de données d'authentification}
    \label{fig:postgresql-auth}
\end{figure}

\begin{figure}[H]
    \centering
    \includegraphics[width=1.0\textwidth]{screenshots/04-databases/postgresql-flight-data.png}
    \caption{Données des vols dans PostgreSQL}
    \label{fig:postgresql-flights}
\end{figure}

\begin{figure}[H]
    \centering
    \includegraphics[width=1.0\textwidth]{screenshots/04-databases/postgresql-staff-data.png}
    \caption{Données du personnel dans PostgreSQL}
    \label{fig:postgresql-staff}
\end{figure}

\subsubsection{MongoDB - Données NoSQL}

\begin{figure}[H]
    \centering
    \includegraphics[width=1.0\textwidth]{screenshots/04-databases/mongodb-collections.png}
    \caption{Collections MongoDB pour le service de chat}
    \label{fig:mongodb-collections}
\end{figure}

\begin{figure}[H]
    \centering
    \includegraphics[width=0.9\textwidth]{screenshots/04-databases/mongodb-messages.png}
    \caption{Documents de messages dans MongoDB}
    \label{fig:mongodb-messages}
\end{figure}

\subsection{Environnement de Développement}

\subsubsection{Structure du Projet}

\begin{figure}[H]
    \centering
    \includegraphics[width=0.8\textwidth]{screenshots/05-development/ide-project-structure.png}
    \caption{Structure complète du projet dans l'IDE}
    \label{fig:project-structure}
\end{figure}

\subsubsection{Code Source Backend}

\begin{figure}[H]
    \centering
    \includegraphics[width=1.0\textwidth]{screenshots/05-development/auth-controller-code.png}
    \caption{Exemple de code - AuthController Spring Boot}
    \label{fig:auth-controller-code}
\end{figure}

\subsubsection{Code Source Frontend}

\begin{figure}[H]
    \centering
    \includegraphics[width=1.0\textwidth]{screenshots/05-development/react-component-code.png}
    \caption{Exemple de code - Composant React TypeScript}
    \label{fig:react-component-code}
\end{figure}

\subsubsection{Tests API avec Postman}

\begin{figure}[H]
    \centering
    \includegraphics[width=1.0\textwidth]{screenshots/05-development/postman-tests.png}
    \caption{Collection de tests API Postman}
    \label{fig:postman-tests}
\end{figure}

\begin{figure}[H]
    \centering
    \includegraphics[width=1.0\textwidth]{screenshots/05-development/postman-auth-test.png}
    \caption{Test d'authentification JWT dans Postman}
    \label{fig:postman-auth-test}
\end{figure}

\subsection{GitHub et Gestion de Version}

\subsubsection{Repository GitHub}

\begin{figure}[H]
    \centering
    \includegraphics[width=1.0\textwidth]{screenshots/06-github/repository-overview.png}
    \caption{Vue d'ensemble du repository GitHub AirOps}
    \label{fig:github-repository}
\end{figure}

\subsubsection{Structure du Code}

\begin{figure}[H]
    \centering
    \includegraphics[width=0.9\textwidth]{screenshots/06-github/code-structure.png}
    \caption{Organisation du code source sur GitHub}
    \label{fig:github-code-structure}
\end{figure}

\subsubsection{Historique des Commits}

\begin{figure}[H]
    \centering
    \includegraphics[width=1.0\textwidth]{screenshots/06-github/commits-history.png}
    \caption{Historique des commits et contributions}
    \label{fig:github-commits}
\end{figure}

\subsection{CI/CD avec GitHub Actions}

\subsubsection{Workflows Configurés}

\begin{figure}[H]
    \centering
    \includegraphics[width=1.0\textwidth]{screenshots/07-cicd/github-actions-workflows.png}
    \caption{Workflows GitHub Actions pour CI/CD}
    \label{fig:github-actions-workflows}
\end{figure}

\subsubsection{Pipeline Backend}

\begin{figure}[H]
    \centering
    \includegraphics[width=1.0\textwidth]{screenshots/07-cicd/backend-pipeline.png}
    \caption{Pipeline de build pour les services backend}
    \label{fig:backend-pipeline}
\end{figure}

\subsubsection{Pipeline Frontend}

\begin{figure}[H]
    \centering
    \includegraphics[width=1.0\textwidth]{screenshots/07-cicd/frontend-pipeline.png}
    \caption{Pipeline de build pour le frontend React}
    \label{fig:frontend-pipeline}
\end{figure}

\subsubsection{Logs de Déploiement}

\begin{figure}[H]
    \centering
    \includegraphics[width=1.0\textwidth]{screenshots/07-cicd/deployment-logs.png}
    \caption{Logs détaillés du processus de déploiement}
    \label{fig:deployment-logs}
\end{figure}

\subsection{Tests et Validation}

\subsubsection{Tests d'Intégration}

\begin{figure}[H]
    \centering
    \includegraphics[width=1.0\textwidth]{screenshots/08-testing/integration-tests.png}
    \caption{Résultats des tests d'intégration automatisés}
    \label{fig:integration-tests}
\end{figure}

\subsubsection{Coverage des Tests}

\begin{figure}[H]
    \centering
    \includegraphics[width=0.9\textwidth]{screenshots/08-testing/test-coverage.png}
    \caption{Couverture de code par les tests unitaires}
    \label{fig:test-coverage}
\end{figure}

\subsection{Métriques et Performance}

\subsubsection{Monitoring des Services}

\begin{figure}[H]
    \centering
    \includegraphics[width=1.0\textwidth]{screenshots/09-monitoring/service-metrics.png}
    \caption{Métriques de performance des microservices}
    \label{fig:service-metrics}
\end{figure}

\subsubsection{Utilisation des Ressources}

\begin{figure}[H]
    \centering
    \includegraphics[width=1.0\textwidth]{screenshots/09-monitoring/resource-usage.png}
    \caption{Utilisation CPU et mémoire en temps réel}
    \label{fig:resource-usage}
\end{figure}

\subsection{Conclusion des Démonstrations}

Ces captures d'écran démontrent la réalisation complète et fonctionnelle de la solution AirOps. L'application offre :

\begin{itemize}
    \item \textbf{Interface utilisateur moderne} : Design responsive avec UX optimisée
    \item \textbf{Architecture robuste} : Microservices scalables et maintenables
    \item \textbf{Sécurité renforcée} : Authentification JWT et contrôle d'accès granulaire
    \item \textbf{Communication temps réel} : WebSocket pour interactions instantanées
    \item \textbf{Persistance polyglotte} : PostgreSQL et MongoDB selon les besoins
    \item \textbf{DevOps moderne} : CI/CD automatisé avec containerisation Docker
    \item \textbf{Monitoring complet} : Surveillance et métriques en temps réel
\end{itemize}

La solution développée répond parfaitement aux exigences fonctionnelles et techniques définies pour la gestion des opérations aéronautiques de l'ONDA, démontrant la viabilité et la qualité de l'architecture microservices mise en place.
